% Combined submission: paper reviews, exercise solution, and presentation preference.
% Group 15 submission.
\documentclass[11pt,a4paper]{article}
\usepackage[T1]{fontenc}
\usepackage[utf8]{inputenc}
\usepackage[english]{babel}
\usepackage{lmodern}
\usepackage[margin=2.2cm]{geometry}
\usepackage{amsmath,amssymb}
\usepackage{booktabs}
\usepackage{microtype}
\usepackage[hidelinks]{hyperref}
\setlength{\parindent}{0pt}
\setlength{\parskip}{0.6em}
\setlength{\emergencystretch}{2em}
\clubpenalty=10000
\widowpenalty=10000
\title{Exercise 1: Paper Reviews and Amdahl's Law}
\author{Michele Arrigoni \and Francesco Foresti \and Kasra Ghasemipoo}
\newcommand{\groupnumber}{15}
\date{Group \groupnumber}

\begin{document}
\maketitle
\vspace{-2em}

\section*{1.1 Paper Reviews}
\subsection*{Evolution of the Graphics Processing Unit (GPU)}
{\small W. J. Dally, S. W. Keckler, and D. B. Kirk.
\emph{IEEE Micro}, vol. 41, no. 6, pp. 42--51, 2021.}

\subsubsection*{1. Primary contribution}
This article provides a historical synthesis of how GPUs evolved from
fixed-function graphics hardware into programmable platforms for graphics,
high-performance computing, and deep learning.
The authors connect major developments, including programmable shaders,
CUDA-enabled computation, and Tensor Cores, to the requirements of these
applications.
By examining hardware alongside programming tools and libraries, they
explain how architectural capabilities and a supporting software ecosystem
jointly enabled the adoption of GPUs beyond graphics.

\subsubsection*{2. Key insight of the contribution}
The key insight is that GPU progress is driven by feedback between applications
and architecture: programmability attracts new workloads, whose requirements
then motivate specialized hardware.
Tensor Cores illustrate this process by accelerating matrix
multiply-and-accumulate operations and amortizing instruction overhead
across many arithmetic operations.
Together with suitable reduced-precision formats and software support,
such specialization explains why the inference-performance trend called
Huang's law exceeds the gains attributable to transistor scaling alone.

\subsubsection*{3. Opinion and reaction}
I found it interesting that the drive for better video games helped create
the computing platform behind the modern AI boom.
The article presents this progression logically and makes the industry's
direction understandable without relying on excessively detailed technical
explanations.
I also appreciated its emphasis on software, which makes clear why
computational power alone is not enough to make a platform useful.

\newpage
\subsection*{A Bridging Model for Parallel Computation}
{\small Leslie G. Valiant.
\emph{Communications of the ACM}, vol. 33, no. 8, pp. 103--111, 1990.}

\subsubsection*{1. Primary contribution}
Valiant introduces the bulk-synchronous parallel (BSP) model as a common
framework between parallel software and hardware, with a role similar to
that of the von Neumann model in sequential computing.
The model combines processing and memory components, point-to-point
communication, and controlled synchronization in units called supersteps.
The paper supports this proposal with results on implementing shared-memory
features and parallel algorithms on BSP, and on implementing BSP efficiently
using different communication networks.

\subsubsection*{2. Key insight of the contribution}
The key insight is that a program should expose more parallel work than
the physical processors can execute at once.
This parallel slackness gives the compiler room to distribute work,
pipeline communication, and absorb uneven memory-access loads.
Together with synchronization at a suitable granularity, it allows many
low-level details to be managed automatically with only a constant-factor
efficiency loss, provided enough parallelism and communication capacity
are available.

\subsubsection*{3. Opinion and reaction}
I found it interesting that the paper treats a shared model for software
and hardware as an essential part of making parallel computing practical.
The idea of exposing more parallelism than the machine can execute at once
seems counterintuitive at first, but its role in handling communication
and uneven workloads makes the argument convincing.
I would have liked a more detailed programming example, since the efficiency
arguments are harder to connect to an implementation than the basic model
itself.

\newpage
\section*{1.2 Amdahl's Law}
\subsection*{1. Derivation}
Let $T_1$ be the sequential execution time of a fixed workload, $P$ the fraction
of this time that can be parallelized, and $N$ the number of processors.
The serial part takes $(1-P)T_1$. Assuming perfect load balancing, no parallel
overhead, and an ideal $N$-fold acceleration of the parallel part, its execution
time is $PT_1/N$. Hence
\begin{equation}
  T_N=(1-P)T_1+\frac{PT_1}{N}
     =T_1\left((1-P)+\frac{P}{N}\right).
\end{equation}
The overall speedup is therefore
\begin{equation}
  \boxed{S(N)=\frac{T_1}{T_N}=\frac{1}{(1-P)+P/N}}.
  \label{eq:law}
\end{equation}
For $0\leq P<1$,
\begin{equation}
  \lim_{N\to\infty}S(N)=\frac{1}{1-P}.
\end{equation}
The serial fraction sets the asymptotic limit. For the special case $P=1$,
the ideal speedup is $S(N)=N$ and there is no finite serial-fraction limit.

\subsection*{2. Correctness for real applications}
Amdahl's law is correct under its assumptions. It is not a universal prediction
for every parallel execution: using $N$ as the acceleration of the parallel
portion assumes that overhead and other architectural effects can be neglected.

\paragraph{Performance predicted too high.}
Communication, synchronization, thread creation, and load imbalance increase
the actual time. Consider a computation with an ideal $100\times$ speedup
whose processors exchange intermediate results
over a network. If communication becomes a bottleneck, the actual speedup
is below $100\times$. Writing
\[
  T_N=T_1\left((1-P)+\frac{P}{N}\right)+T_{\mathrm{overhead}}(N)
\]
makes clear why positive overhead reduces the ideal speedup.

\paragraph{Performance predicted too low.}
A working set that exceeds the cache of one processor may fit into the aggregate caches
when partitioned across several processors. Reducing slow memory accesses
can accelerate the parallel part by more than $N$, so actual speedup can
exceed the prediction of Equation~\eqref{eq:law}, and can even be superlinear.
This changes the assumed acceleration of the parallel portion; it does not
invalidate the time-decomposition argument.

\paragraph{Changing the workload.}
In weather forecasting, more processors are often used to increase grid
resolution rather than to finish an unchanged
problem sooner. Amdahl's fixed-workload limit does not directly characterize
this comparison. Weak scaling studies increasing the workload in proportion
to the number of processors, while Gustafson's law addresses scaled speedup
using the serial fraction of a parallel run. These perspectives change the
workload assumptions used in the fixed-size Amdahl model.

\subsection*{3. Numerical calculations and comparison}
\subsubsection*{(a) $P=0.5$ and $N=4$}
\[
  S_A=\frac{1}{(1-0.5)+0.5/4}
     =\frac{1}{0.5+0.125}
     =\boxed{1.6}.
\]
The parallel execution takes $0.625T_1$, and its efficiency is
$S_A/4=0.4=40\%$.

\subsubsection*{(b) $P=0.2$ and $N=8$}
\[
  S_B=\frac{1}{(1-0.2)+0.2/8}
     =\frac{1}{0.8+0.025}
     =\boxed{\frac{40}{33}\approx1.212}.
\]
The parallel execution takes $0.825T_1$, and its efficiency is
$S_B/8=5/33\approx15.15\%$.

\subsubsection*{(c) Comparison}
\begin{center}
\begin{tabular}{@{}lrrrr@{}}
\toprule
Case & $P$ & $N$ & Speedup & Limit as $N\to\infty$\\
\midrule
A & $0.5$ & $4$ & $1.6$ & $2$\\
B & $0.2$ & $8$ & $1.212\ldots$ & $1.25$\\
\bottomrule
\end{tabular}
\end{center}
Case A is faster relative to its sequential baseline despite having fewer
processors. Its larger parallelizable fraction leaves a smaller serial
bottleneck. In case B, $80\%$ of the original time remains serial, so even
arbitrarily many processors cannot produce a speedup above $1.25$.

These cases change both $P$ and $N$, so their difference cannot be attributed
to processor count alone. For any fixed $P>0$, adding processors increases
ideal speedup, but with diminishing returns when $P<1$.
Reducing the serial fraction is essential for effective scaling.

For comparison, keeping $P=0.5$ and increasing only $N$ from four to eight
would give
\[
  S(8)=\frac{1}{0.5+0.5/8}=\frac{16}{9}\approx1.778.
\]


\section*{1.3 Willingness to Present}
We are willing to present the following in English:
\begin{itemize}
  \item Paper review 1: \emph{Evolution of the Graphics Processing Unit (GPU)}.
  \item Paper review 2: \emph{A Bridging Model for Parallel Computation}.
  \item The exercise on Amdahl's law.
\end{itemize}

\end{document}
